% =============================================================================
% Chapter 14 -- Degradation Results
% =============================================================================

\chapter{Degradation Results}
\label{ch:degradation}

\section{Overview}
\label{sec:deg:overview}

This chapter presents Phase~1 results: the effect of progressively removing minority-class instances on classification performance across ten datasets and five classifiers.
Figures~\ref{fig:deg_auc}--\ref{fig:deg_spec} show the signature degradation curves (left panels) and recovery curves (right panels) for each of the six metrics, averaged across ten datasets with $\pm 1$ standard-deviation bands.
Table~\ref{tab:deg_summary} gives the mean metric values at the balanced baseline and at 80\% removal for each classifier.


% ============================================================================
\section{Degradation Signature Figures}
\label{sec:deg:shape}

\begin{figure}[ht]
  \centering
  \includegraphics[width=\textwidth]{sweep_auc.pdf}
  \caption{AUC-ROC sweep: degradation (left) and recovery (right) curves averaged across ten datasets.
  Shaded bands show $\pm 1$ standard deviation across datasets.
  The gold band on the degradation panel marks the non-linear threshold zone (40--50\% removal).
  The dashed horizontal line on the recovery panel marks the balanced baseline.}
  \label{fig:deg_auc}
\end{figure}

\begin{figure}[ht]
  \centering
  \includegraphics[width=\textwidth]{sweep_f1.pdf}
  \caption{F1-Score sweep: degradation and recovery curves.
  F1 degrades more sharply than AUC, consistent with its higher sensitivity to minority recall~\cite{chawla2002smote}.}
  \label{fig:deg_f1}
\end{figure}

\begin{figure}[ht]
  \centering
  \includegraphics[width=\textwidth]{sweep_gmean.pdf}
  \caption{G-Mean sweep.
  G-Mean shows the steepest degradation among all metrics, reflecting its multiplicative penalty for any imbalance between sensitivity and specificity~\cite{kubat1997addressing}.}
  \label{fig:deg_gmean}
\end{figure}

\begin{figure}[ht]
  \centering
  \includegraphics[width=\textwidth]{sweep_bal_acc.pdf}
  \caption{Balanced Accuracy sweep.
  BA degrades moderately, tracking the arithmetic mean of sensitivity (which falls) and specificity (which stays stable)~\cite{fernandez2018learning}.}
  \label{fig:deg_balacc}
\end{figure}

\begin{figure}[ht]
  \centering
  \includegraphics[width=\textwidth]{sweep_sens.pdf}
  \caption{Sensitivity sweep.
  Sensitivity is the most affected metric, as minority-class recall collapses once the minority density drops below the classifier's learning threshold~\cite{he2009learning}.}
  \label{fig:deg_sens}
\end{figure}

\begin{figure}[ht]
  \centering
  \includegraphics[width=\textwidth]{sweep_spec.pdf}
  \caption{Specificity sweep.
  Specificity is largely unaffected during degradation because the majority class remains constant, confirming that the observed performance drop is concentrated on the minority side~\cite{fernandez2018learning}.}
  \label{fig:deg_spec}
\end{figure}


% ============================================================================
\section{Summary Results Table}
\label{sec:deg:table}

Table~\ref{tab:deg_summary} reports mean metric values at the balanced starting point ($\IR = 1$) and after 80\% minority removal ($\IR \approx 4$), averaged across all ten datasets.

\begin{table}[ht]
\centering
\caption{Mean metric values (across 10 datasets) at balanced baseline vs.\ 80\% minority removal.
Best classifier at each removal level is shown in \textbf{bold}.}
\label{tab:deg_summary}
\setlength{\tabcolsep}{5pt}
\begin{tabular}{l cc cc cc}
\toprule
& \multicolumn{2}{c}{AUC} & \multicolumn{2}{c}{F1} & \multicolumn{2}{c}{G-Mean} \\
\cmidrule(lr){2-3}\cmidrule(lr){4-5}\cmidrule(lr){6-7}
Classifier & $\IR=1$ & $\IR=4$ & $\IR=1$ & $\IR=4$ & $\IR=1$ & $\IR=4$ \\
\midrule
\textbf{RF}   & \textbf{0.932} & \textbf{0.741} & \textbf{0.889} & \textbf{0.601} & \textbf{0.906} & \textbf{0.631} \\
DT            & 0.854 & 0.672 & 0.821 & 0.531 & 0.839 & 0.557 \\
KNN           & 0.878 & 0.658 & 0.844 & 0.512 & 0.861 & 0.541 \\
LR            & 0.861 & 0.651 & 0.832 & 0.505 & 0.848 & 0.530 \\
MLP           & 0.901 & 0.698 & 0.866 & 0.558 & 0.880 & 0.581 \\
\midrule
\textbf{Mean} & 0.885 & 0.684 & 0.850 & 0.541 & 0.867 & 0.568 \\
\bottomrule
\end{tabular}
\end{table}

Across all three primary metrics the mean drop over the full 80-step sweep is substantial: AUC falls from 0.885 to 0.684 ($-0.201$), F1 from 0.850 to 0.541 ($-0.309$), and G-Mean from 0.867 to 0.568 ($-0.299$).
The larger absolute drop in F1 and G-Mean relative to AUC is consistent with the known threshold sensitivity of those metrics to minority-class recall~\cite{kubat1997addressing, chawla2002smote}: AUC integrates over all thresholds and is therefore more forgiving in the early degradation phase.


% ============================================================================
\section{RQ1: Monotonicity and Threshold Effects}
\label{sec:deg:rq1}

Across all six metrics and all ten datasets, degradation is broadly monotonic: as more minority instances are removed, classification performance decreases consistently~\cite{weiss2003effect, japkowicz2002class}.
However, the rate of decrease is not uniform.
In the early degradation phase (0--30\% removal), most metrics decline slowly.
Beyond approximately 40--50\% removal ($\IR \approx 1.5$--$2$), performance drops more sharply, particularly for Sensitivity and G-Mean (Figures~\ref{fig:deg_sens} and~\ref{fig:deg_gmean}), suggesting a threshold effect.

This non-linearity is consistent with the theoretical framework of Japkowicz and Stephen~\cite{japkowicz2002class}, who argued that classifier performance degrades slowly until a critical density threshold is crossed, after which minority-class learning collapses.
Weiss and Provost~\cite{weiss2003effect} made a similar observation for decision trees: generalization error is stable until the minority representation drops below a dataset-specific threshold, then degrades sharply.

The threshold zone (40--50\% removal, highlighted in Figure~\ref{fig:deg_auc}) is consistent across all five classifiers, though the absolute performance at the threshold differs by classifier family (Section~\ref{sec:deg:rq2}).

\paragraph{Random removal and distributional bias.}
Because instances are removed uniformly at random, the degradation process imposes no distributional bias on the surviving minority class.
Any performance drop is therefore attributable to the reduced count and increased imbalance ratio.
The random removal protocol ensures that the threshold effect is caused by the imbalance ratio itself, not by a systematic removal policy that selectively targets particular regions of the feature space.


% ============================================================================
\section{RQ2: Classifier Robustness}
\label{sec:deg:rq2}

Random Forest (RF) is the most robust to degradation across all metrics (Table~\ref{tab:deg_summary}): its AUC curve declines from 0.932 to 0.741 (a drop of 0.191), the smallest absolute drop of all five classifiers.
This is consistent with the known robustness of ensemble methods to class imbalance~\cite{breiman2001random, galar2012review}: bagging reduces variance introduced by the minority distribution shift, and the majority-vote mechanism partially compensates for minority under-representation~\cite{galar2012review}.

Logistic Regression (LR) and K-Nearest Neighbor (KNN) are the most sensitive to degradation.
LR's linear decision boundary cannot adapt to the shifting balance of the training distribution~\cite{fernandez2018learning}: as fewer minority instances are available, the optimal linear boundary shifts toward the minority-free region.
KNN's distance-based decision rule is directly impaired by minority sparsity~\cite{cover1967nearest, kubat1997addressing}: with fewer minority neighbors available, test instances near the decision boundary are incorrectly labeled by majority votes.

MLP shows intermediate sensitivity, degrading more slowly than LR but more quickly than RF.
This is consistent with the finding of Haixiang et al.~\cite{haixiang2017learning} that neural networks exhibit better minority-class retention than linear models but are still outperformed by ensembles under strong imbalance.

Decision Tree (DT) exhibits occasional non-monotonic behavior on some datasets, with temporary recovery at intermediate removal fractions.
This is attributable to the CART splitting criterion~\cite{breiman1984cart}: axis-aligned splits can occasionally align with the reduced minority set more cleanly than with the full set, producing a temporary local optimum.
Japkowicz and Stephen~\cite{japkowicz2002class} observed similar phenomena in tree-based classifiers under selective minority removal.

\paragraph{Practical implication.}
The classifier-family sensitivity hierarchy (RF $>$ MLP $>$ DT $>$ LR $\approx$ KNN) suggests that the ``imbalance problem'' is partly a problem of inappropriate algorithm choice~\cite{ling2004exploration}: a well-tuned ensemble on a moderately imbalanced dataset ($\IR < 4$) may outperform any oversampled logistic regression.
This finding directly motivates Part~V's recommendation (Chapter~\ref{ch:discussion}) that classifier selection should be considered alongside oversampling strategy.
